\documentclass[10pt,a4paper]{article}
\usepackage[T1]{fontenc}
\usepackage[utf8]{inputenc}
\usepackage{lmodern,amsmath,amssymb,booktabs,array,microtype}
\pdfmapfile{+lm.map}
\pdfmapfile{+cm.map}
\pdfmapfile{+cmextra.map}
\pdfmapfile{+symbols.map}
\pdfinfoomitdate=1
\pdftrailerid{}
\usepackage[a4paper,margin=23mm]{geometry}
\usepackage{xcolor,hyperref,fancyhdr}
\definecolor{ink}{HTML}{183A4A}
\hypersetup{colorlinks=true,linkcolor=ink,urlcolor=ink,pdftitle={CST-DUST: Definitionsbereich des gekoppelten Kerns v0.1}}
\pagestyle{fancy}\fancyhf{}\fancyhead[L]{\small CST-DUST | Gekoppelter Kern}\fancyhead[R]{\small v0.1 | 11.11.2026}\fancyfoot[C]{\thepage}
\setlength{\headheight}{14pt}\setlength{\parindent}{0pt}\setlength{\parskip}{5pt}
\newcommand{\HH}{\mathcal H_A}\newcommand{\VV}{\mathcal V}
\newcommand{\dd}{\,\mathrm d}\newcommand{\ii}{\mathrm i}
\newcommand{\R}{\mathbb R}\newcommand{\norm}[1]{\left\lVert#1\right\rVert}
\newcommand{\abs}[1]{\left|#1\right|}
\newcommand{\one}{\mathbf1}\newcommand{\Ccal}{\mathcal C_{33}}
\begin{document}
{\color{ink}\LARGE\bfseries Definitionsbereich des gekoppelten Kerns}\par
{\large Exaktes Gegenbeispiel zur $L^2$-Evolution und kontrollierte Testzustandspaarungen}\par
\textbf{Version 0.1, 11. November 2026.} Anschluss an PR \#279--283.\\
Projekt: \texttt{thomasmgfeller/x6-spectral-astrolabium}.\\
Geprüfte Basis:
\texttt{cfd4413c8e9d2338f41567aac249e9ac54a295d8}.

\section*{Ergebnis}
Die \emph{ursprüngliche gekoppelte} AA-Kontrollamplitude bildet den bereits
verwendeten normierten Zustand $\psi=\Psi_{0,00}$ nicht in den gewählten
Rand-Hilbertraum ab:
\[
 \norm{\psi}_{\HH}=1,\qquad K_0\psi\notin\HH.
\]
Der Beweis berücksichtigt alle 33 Historien, ihre unveränderten komplexen
Gravitationsphasen, die Randquotienten und beide inneren Dichtestrata.
Er benutzt weder die entkoppelte Referenz noch die Divergenz eines isolierten
Historienbeitrags. Auf dem ausgehenden Atomsektor verbleiben zwei unterschiedliche
quadratische Feldphasen. Ihr Kreuzterm kann das positive, linear wachsende
Normquadrat nicht auslöschen.

Gleichzeitig ist derselbe Kern ohne Änderung von Wirkung oder Maß als beschränkte
Abbildung $K_0:\VV\to\VV'$ wohldefiniert. Dabei
\[
 \VV=L^2_{\rm inv}\bigl((1+|q|^2)^4\dd\mu_A\bigr)
 \subset\HH\subset
 \VV'=L^2_{\rm inv}\bigl((1+|q|^2)^{-4}\dd\mu_A\bigr).
\]
Das Strichsymbol bezeichnet den Antidualraum bezüglich der unveränderten
sesquilinearen Randpaarung. $\VV'$ ist ein Raum verallgemeinerter Antworten,
\emph{kein neu postuliertes physikalisches Skalarprodukt}.

\begin{center}\begin{tabular}{p{.40\linewidth}p{.51\linewidth}}\toprule
Frage & Ergebnis im deklarierten endlichen Modell\\\midrule
Normierter Eingang ergibt $L^2$-Ausgang? & Gegenbeispiel $\Psi_{0,00}$; Antwortnorm unendlich.\\
Können sich alle problematischen Beiträge auslöschen? & Für diesen Eingang nein; vollständige 33er-Summe geprüft.\\
Sind die bisherigen Paketpaarungen gültig? & Ja, als Paarungen $\VV\times\VV$ mit $K_0:\VV\to\VV'$.\\
Ist der Abel-Grenzwert kontrolliert? & Ja, Konvergenz sogar in $\VV\to\VV'$-Operatornorm, Fehler $<1.342\epsilon$.\\
Physikalischer CST-Evolutionskern? & Weiter \texttt{NOT\_CONSTRUCTED}; alle acht Gates \texttt{HOLD}.\\\bottomrule
\end{tabular}\end{center}

\textbf{Geltungsbereich.} Vollständig ist hier die bereits deklarierte endliche
Historienfamilie. Eine unendliche oder vierdimensionale physikalische CST-Summe
ist damit weder definiert noch ausgewertet. Der Satz widerlegt die Evolutionsdeutung
dieses Kontrollkerns im bisherigen Hilbertraum; er ist kein allgemeiner
Ausschlusssatz für CST.
\newpage

\section{Unveränderter Kern, Zustände und Maße}
Es gelten $L=\{0,1\}$, $I=\{2\}$, $R=\{3,4\}$, die 33 abgedeckten AA-Ordnungen
und die Kontrollwerte $\kappa=1$, $\mu=0$, $\lambda_{\rm Kopplung}=1$.
Mit den bisherigen dimensionslosen Variablen und $\beta=4/\sqrt6$:
\[
 \dd\lambda(t)=\tfrac12\one_{[-1,1]}\dd t,\qquad
 \dd\nu(r)=\tfrac12\delta_0(\dd r)+\tfrac12\one_{(0,1)}\dd r,
 \qquad
 \dd\mu_A=\tfrac12\prod_{j\in A}\dd\lambda(t_j)\dd\nu(r_j)\dd q_j.
\]
Die zwei Randpunkte werden mit ihrer gesamten Ausstattung gemeinsam vertauscht;
$\HH$ ist der invariante Unterraum von $L^2(\dd\mu_A)$.
Die zeit- und dichteabhängige Maßkomponente $\dd\eta_A$ hat Gesamtmasse $1/2$;
die Feldkomponente ist das ungegewichtete Lebesgue-Maß auf $\R^2$.

Für jede Ordnung sei $c_{ij}\in\{1,-9,16,-8\}$ der unveränderte
Intervallkoeffizient und
\[
 \frac{S_d}{\hbar}=-\sum_{j=0}^4 r_j k_j(t),\quad
 k_j(t)=\frac12+\frac\beta4\sum_{i\prec j}c_{ij}(t_j-t_i)^2,
 \qquad F_4(C)=5-\sum_{i\prec j}c_{ij}.
\]
Die skalare Hesse-Matrix ist $M_C=\beta H_C^{(0)}$, wobei jedes Paar $(i,j)$
die Matrixeinträge $H_{ii},H_{jj}\mathrel{-}=c_{ij}/2$ und
$H_{ij},H_{ji}\mathrel{+}=c_{ij}/2$ beiträgt. In der Reihenfolge
$b=(q_0,q_1,q_3,q_4)$, $q_2$ schreiben wir
\[
 M_C=\begin{pmatrix}B_C&J_C^T\\J_C&-\alpha_C\end{pmatrix},\quad
 \beta\le\alpha_C\le2\beta,\quad \norm{J_C}^2\le\beta^2.
\]
Die Lücken und Normschranken werden für alle 33 Matrizen exakt geprüft.
Die innere Feldintegration mit $\epsilon\ge0$ lautet
\[
 G_{C,\epsilon}(b)=\frac{\sqrt{2\pi}}{\sqrt{\epsilon+\ii\alpha_C}}
 \exp\left(\frac\ii2b^TB_Cb-
          \frac{(J_Cb)^2}{2(\epsilon+\ii\alpha_C)}\right).
\]
Die Wurzel wird aus der rechten Halbebene fortgesetzt; bei null entsteht die
Fresnelphase $e^{-\ii\pi/4}$. Mit $z_L=(t_0,r_0,t_1,r_1)$ und entsprechend $z_R$:
\[
 K_\epsilon(z_R,q_R;z_L,q_L)=\frac1{33}\sum_{C\in\Ccal}e^{\ii F_4(C)}
 \int\dd\lambda(t_2)\dd\nu(r_2)\,
 e^{-\ii\sum_j r_jk_j(t)}G_{C,\epsilon}(b).
 \tag{1}
\]
Es wird weder durch eine komplexe Zustandssumme geteilt noch ein weiterer
Orbit- oder Automorphismusfaktor eingeführt. Der Kern ist invariant unter beiden
Randvertauschungen, da diese die 33 Summanden nur permutieren.

Die innere Integration ist ein Abel/Fresnel-Integral, kein absolut konvergentes
unreguliertes Feldintegral. Für reelle Randfelder ist aber
\[
 |K_\epsilon|\le C_0:=\sqrt{2\pi/\beta}<1.962,
 \qquad \epsilon\ge0. \tag{2}
\]
Dies folgt aus $|G_{C,\epsilon}|\le\sqrt{2\pi/\alpha_C}$ und der normierten
inneren Staubmaßkomponente. Damit ist $K_0$ eine konkret definierte beschränkte
Kernfunktion; daraus folgt noch keine Beschränktheit als $L^2$-Operator.
\newpage

\section{Testraum, Definitionsbereich und Fehlerbegriff}
Für $f\in\VV$ sei
\[
 \norm f_\VV^2=\int(1+|q|^2)^4|f(z,q)|^2\dd\mu_A.
\]
Dieser Raum ist dicht in $\HH$: Feldabschneidungen und symmetrisierte kompakt
getragene einfache Funktionen approximieren jedes $L^2$-Element. Für alle
$f,g\in\VV$ existiert die Paarung
\[
 \mathcal B_\epsilon(g,f)=\iint\overline{g(y)}K_\epsilon(y,x)f(x)
 \dd\mu_A(y)\dd\mu_A(x).
 \tag{3}
\]
Nach der inneren Abel-Integration ist dieses Randintegral absolut konvergent.
In zwei Felddimensionen gelten exakt
\[
 \int\frac{\dd\mu_A}{(1+|q|^2)^4}=\frac\pi6,
 \qquad
 \int\frac{|q|^4\dd\mu_A}{(1+|q|^2)^4}=\frac\pi6.
 \tag{4}
\]
Cauchy--Schwarz liefert $\norm f_{L^1}\le\sqrt{\pi/6}\norm f_\VV$ sowie
$\norm{|q|^2f}_{L^1}\le\sqrt{\pi/6}\norm f_\VV$. Deshalb
\[
 |\mathcal B_\epsilon(g,f)|\le C_4\norm g_\VV\norm f_\VV,
 \quad \norm{K_\epsilon}_{\VV\to\VV'}\le C_4=\frac\pi6C_0<1.028.
 \tag{5}
\]
Dies legt die zulässige Interpretation als kontinuierliche sesquilineare
Amplitude fest. Hermite-Gaußpakete mit endlichem Polynomgrad sowie die bisherigen
Box-Uhr- und Maskenpakete gehören zu $\VV$; für $\psi=\Psi_{0,00}$ gilt
$\norm\psi_\VV^2=65$, während $\norm\psi_{\HH}=1$ bleibt.

\textbf{Abel-Fehler.} Differenzieren des exakt integrierten $G$ ergibt
\[
 |\partial_\epsilon G_{C,\epsilon}(b)|
 \le C_0\left(\frac1{2\beta}+\frac{|b|^2}{2}\right).
\]
Wegen (4), $|b|^2=|q_L|^2+|q_R|^2$ und des positiven normierten inneren Maßes:
\[
 \norm{K_\epsilon-K_0}_{\VV\to\VV'}
 \le \epsilon C_4\left(1+\frac1{2\beta}\right)
 <1.342\epsilon. \tag{6}
\]
Es handelt sich um eine gleichmäßige Normschranke, nicht nur um Konvergenz
endlich vieler Matrixeinträge. Spezifische frühere Paketfehler können enger sein.

\textbf{Feldabschneidung.} Wird (3) auf $|q_L|\le R_L$, $|q_R|\le R_R$
beschränkt, ist der Fehler höchstens
\[
 C_4\big[(1+R_L^2)^{-3/2}+(1+R_R^2)^{-3/2}\big]
 \norm g_\VV\norm f_\VV. \tag{7}
\]
Denn das gewichtete äußere Maß ist $\pi(1+R^2)^{-3}/6$.
Ein Fehler $|\widetilde K-K|\le\delta$ auf einem Integrationsbereich ergibt
dort zusätzlich $\delta\pi\norm g_\VV\norm f_\VV/6$.
Paketapproximationen erfüllen
\[
 |\mathcal B(g,f)-\mathcal B(g_N,f_N)|\le C_4\left[
 \norm{g-g_N}_\VV\norm f_\VV+
 \norm{g_N}_\VV\norm{f-f_N}_\VV\right]. \tag{8}
\]
Konvergenz nur in der ungegewichteten Hilbertnorm genügt dafür nicht.
Numerische Quadratur- und Rundungsfehler sind zusätzlich zu zertifizieren.
\newpage

\textbf{Exakter Hilbert-Definitionsbereich innerhalb des Testraums.}
Ohne Erweiterungsannahme ist
\[
 D_{\HH}(K_0)=\{f\in\VV:K_0f\in\HH\},\qquad
 K_{\HH}=K_0|_{D_{\HH}(K_0)}. \tag{9}
\]
Äquivalent, für eine orthonormale Basis $(e_j)$ von $\HH$, deren endlicher
linearer Spann dicht in $\VV$ ist:
\[
 f\in D_{\HH}(K_0)\quad\Longleftrightarrow\quad
 \sup_N\sum_{j=1}^N|\mathcal B_0(e_j,f)|^2<\infty. \tag{10}
\]
Eine solche Basis entsteht durch Gram--Schmidt bezüglich $\HH$ aus einer
abzählbaren in $\VV$ dichten Testfamilie. Der Raum $\VV$ ist separabel;
das Weglassen endlicher linearer Abhängigkeiten verändert den linearen Spann
nicht. Weil $\VV$ dicht in $\HH$ ist, ist diese Basis auch dort vollständig.
Die Notwendigkeit von (10) folgt aus Parseval. Für die Umkehr definieren die
quadratsummierbaren
Koeffizienten definieren dann einen $\HH$-Vektor, der auf einer $\VV$-dichten
Familie mit dem stetigen Funktional übereinstimmt und somit überall.
\emph{Die Totalität in $\VV$ ist Teil dieser Basiswahl; beliebige nur in $\HH$
vollständige Tests reichen für diesen Umkehrschluss nicht aus.}

(9)--(10) sind eine präzise Charakterisierung des maximalen $\HH$-wertigen
Bereichs \emph{innerhalb von $\VV$}. Eine vollständige Klassifikation aller
Elemente dieses Bereichs, seine Dichtheit in $\HH$, Abgeschlossenheit oder
Erweiterung auf Eingänge außerhalb von $\VV$ wird nicht behauptet.
Für die angefragte Evolutionsdeutung reicht das folgende konkrete Gegenbeispiel.

\section{Normierter Zeuge und vollständige Historienklassifikation}
Wir verwenden unverändert
\[
 \psi(z_L,x)=2\sqrt2\one_{r_0=r_1=0}\,h_0(x),\qquad
 h_0(x)=\pi^{-1/2}e^{-|x|^2/2}.
 \tag{11}
\]
Die Eingangsuhrfunktion ist eins. Das Dichteatom und der Randquotient geben
$\norm\psi_{\HH}^2=(1/2)(1/4)(2\sqrt2)^2=1$.
Untersucht wird die \emph{gesamte Ausgangsantwort} auf dem Atomsektor
$r_3=r_4=0$. Dieser ist keine Nullmenge: seine Randmaßkomponente hat Masse $1/8$.
Andere ausgehende Dichtestrata können Beiträge an diesen Punkten nicht auslöschen.

Sei $\widehat K_C=H_{BB}^{(0)}-H_{B2}^{(0)}(H_{22}^{(0)})^{-1}H_{2B}^{(0)}$
das rationale Schur-Komplement, mit Eingangsblock $L_C$, Kreuzblock $X_C$
und Ausgangsblock $R_C$. Die Feldantwort nach Eingangsintegration hat Form
\[
 I_C(y)=\frac{\sqrt{2\pi/\alpha_C}\,e^{-\ii\pi/4}\,2\sqrt\pi}
                 {\sqrt{\det_{\rm cont}(I-\ii\beta L_C)}}
 \exp\left[\frac{\ii\beta}2y^TR_Cy-
 \frac{\beta^2}2y^TX_C^T(I-\ii\beta L_C)^{-1}X_Cy\right]. \tag{12}
\]
Die kontinuierliche Determinantenwurzel ist das Produkt der Hauptwurzeln
$\sqrt{1-\ii\beta\ell_j}$ der reellen Eigenwerte $\ell_j$ von $L_C$.
Es wird keine unkontrollierte Hauptwurzel des Gesamtprodukts verwendet.
Der reelle Dämpfungsblock ist
\[
 Q_C=\beta^2 X_C^T(I+\beta^2L_C^2)^{-1}X_C\succeq0,
 \qquad\ker Q_C=\ker X_C. \tag{13}
\]
\newpage

Alle 33 Kreuzblöcke werden exakt, vor einer numerischen Näherung, klassifiziert:
\begin{center}\begin{tabular}{llll}\toprule
Historien & Anzahl & Rang $X_C$ & Ausgehende Nullrichtung in $(q_3,q_4)$\\\midrule
Alle übrigen & 28 & 2 & keine\\
9116 & 1 & 1 & $(3/26,1)$\\
17308 & 1 & 1 & $(26/3,1)$\\
25372, 25496 & 2 & 1 & $(-1,1)$\\
25500 & 1 & 1 & $(-1,1)$\\\bottomrule
\end{tabular}\end{center}
Für $s=(q_3+q_4)/\sqrt2$, $a=(q_3-q_4)/\sqrt2$ und beschränktes $s$ zerfallen
somit 30 Beiträge gaußförmig für $|a|\to\infty$. Die beiden schiefen Nullrichtungen
liegen nicht entlang der betrachteten $a$-Achse. Nur die letzten drei Beiträge
bleiben auf diesem Streifen ungedämpft.

\section{Staubsumme und die zwei verbleibenden Phasen}
Auf dem ausgehenden Atomsektor und für (11) liegen alle äußeren Dichten auf null.
Das bedeutet \emph{nicht}, dass die innere Dichte auf null gesetzt wird. Es bleibt
der vollständige Faktor
\[
 D_p=\int\dd\lambda(t_0)\dd\lambda(t_1)\dd\lambda(t_2)\,
 \frac{1+E(k_2)}2,\quad
 E(k)=\int_0^1e^{-\ii kr}\dd r,
 \quad k_2=\frac12+\frac\beta4\sum_{i\prec2}(t_2-t_i)^2. \tag{14}
\]
Dabei ist $p\in\{1,2\}$ die Vorgängerzahl von Punkt 2; die unbenutzte
Eingangsuhr integriert zu eins. Die Vorgänger sind stets direkte Nachbarn.
Für beide Fälle gilt $0<k_2\le1/2+2\beta<4$. Auf $0\le r\le1/4$ ist daher
$\cos(k_2r)\ge1-(k_2r)^2/2\ge1/2$, und auf dem Rest ist $1+\cos(k_2r)\ge0$.
Somit, ohne Quadratur oder Annahme über kleine Phasen:
\[
 \Re D_p\ge\frac12\int_0^{1/4}\frac32\dd r=\frac3{16}>0,
 \qquad |D_p|\le1. \tag{15}
\]
Beide inneren Masken und alle Eingangs-/Innenuhren sind also berücksichtigt.

Für 25372 gilt
\[
 L_1=\begin{pmatrix}26/3&0\\0&-1\end{pmatrix},\quad
 X_1=\begin{pmatrix}-13/3&-13/3\\1/2&1/2\end{pmatrix},\quad
 R_1=\begin{pmatrix}11/3&1/6\\1/6&11/3\end{pmatrix},\quad\alpha_1=3\beta/2.
\]
25496 ist die Eingangsvertauschung davon. Bei symmetrischem Eingang und (14)
sind beide Antworten exakt gleich, einschließlich $D_1$ und $e^{18\ii}$.
Für 25500 gilt
\[
 L_2=R_2=\begin{pmatrix}69/8&1/8\\1/8&69/8\end{pmatrix},\quad
 X_2=-\frac{35}8\begin{pmatrix}1&1\\1&1\end{pmatrix},\quad\alpha_2=2\beta.
\]
Seine Gravitationsphase ist $e^{37\ii}$, sein Staubfaktor $D_2$.
Definiere $f_j(s)$ als (12) für $y=(s,s)/\sqrt2$ der jeweiligen Matrix.
Dann ist die vollständige Antwort auf diesem ausgehenden Atomsektor
\[
 g(s,a)=L_*\left[2e^{18\ii}D_1f_1(s)e^{\ii\omega_1 a^2}
                  +e^{37\ii}D_2f_2(s)e^{\ii\omega_2 a^2}\right]
                  +\mathcal R(s,a),\quad
 L_*=\frac{\sqrt2}{132},\quad
 (\omega_1,\omega_2)=\frac\beta4(7,17). \tag{16}
\]
$L_*$ enthält den Eingangsquotienten, die Eingangsatommasse, die normierte
Paketamplitude und $1/33$. Die restlichen 30 Historien bilden $\mathcal R$;
sie wurden nicht aus der Summe entfernt.
\newpage

\section{Keine Auslöschung: positiver Grenzwert der Normdichte}
Auf $|s|\le1$ gilt für hinreichend großes $|a|$ gleichmäßig
$|\mathcal R(s,a)|\le A e^{-c a^2}$ mit gewissen $A,c>0$.
Dies folgt für jeden der 30 Summanden aus (13), $X_C(1,-1)^T\ne0$
und der Endlichkeit der Familie. Ihre Beiträge sowie ihre Kreuzterme mit den
beschränkten persistierenden Summanden integrieren sich zu $O(1)$.

Auch der Kreuzterm zwischen den beiden persistierenden Frequenzen ist nur
$O(1)$: $\omega_2-\omega_1=5\beta/2\ne0$ und
$\int_{-R}^R e^{\ii(\omega_2-\omega_1)a^2}\dd a$ ist gleichmäßig beschränkt.
Zum Beispiel liefert partielle Integration außerhalb $|a|\le1$ eine konstante
Schranke; die Ableitung der Phase ist dort von null getrennt.
Die beiden komplexen Koeffizienten und ihre Gravitationsphasen ändern daran nichts.

Sei $\Omega_R$ der Ausgangsbereich $r_3=r_4=0$, $|s|\le1$, $|a|\le R$
mit allen Ausgangsuhren. Dann folgt exakt aus (16):
\[
 \lim_{R\to\infty}\frac{\norm{\one_{\Omega_R}g}_{\HH}^2}{2R}
 =\frac{L_*^2}{8}\int_{-1}^{1}
       \left(4|D_1f_1(s)|^2+|D_2f_2(s)|^2\right)\dd s>0. \tag{17}
\]
Der Faktor $1/8$ ist die ausgehende Atommaßkomponente einschließlich Quotient.
Die Feldtransformation hat Jacobi-Betrag eins; die Ausgangsuhren haben Maß eins.
Es gibt keinen negativen Normbeitrag aus den anderen Dichtestrata.

\textbf{Explizite untere Schranke.} Exakte Gaußalgebra ergibt
\[
 |f_2(s)|^2=P_2\,e^{-c_2s^2},\quad
 P_2=\frac{8\pi^2}{2\beta\sqrt{(1231/6)(581/3)}},\quad
 c_2=\frac{1225}{1231}<1. \tag{18}
\]
Die entsprechende Dämpfung von $|f_1|^2$ ist $51484/59785<1$.
Schon das Weglassen des positiven ersten Diagonalterms in (17) liefert
\[
 \lim_{R\to\infty}\frac{\norm{\one_{\Omega_R}g}_{\HH}^2}{2R}
 \ge\frac1{8712\cdot8}\left(\frac3{16}\right)^2\frac{2P_2}{e}
 =\frac{3\sqrt{2145633}\,\pi^2}{354470014976\,e}
 >4.5\cdot10^{-8}. \tag{19}
\]
Der Wert der ausgewerteten unteren Schranke ist
$4.5011629414369758512\ldots\cdot10^{-8}$; die letzte strikte Ungleichung
wird zusätzlich mit 192-Bit-Ballarithmetik eingeschlossen.
Die Zahl ist eine \emph{asymptotische untere Steigung}, kein bereits für jedes
endliche $R$ behaupteter Schnittfehler oder physikalischer Messwert.

(17)--(19) beweisen $\norm g_{\HH}=\infty$ und damit
\[
 \boxed{\Psi_{0,00}\notin D_{\HH}(K_0),\qquad K_0\Psi_{0,00}\in\VV'\setminus\HH.}
\]
Das ist der verlangte Gegenbeleg für die ursprüngliche gekoppelte Amplitude.
Für jeden endlichdimensionalen Hilbertvektor $g_N$ gilt außerdem
$g-g_N\notin\HH$; andernfalls wäre auch $g$ als Summe zweier Hilbertvektoren
quadratintegrierbar. Ein endlicher Satz genauer Matrixelemente kann diesen
Definitionsbereichsfehler daher nicht durch eine globale $L^2$-Restschranke beheben.
\newpage

\section{Nachrechnung, zulässige Aussagen und offener Anschluss}
Das Prüfpaket enthält die unveränderten Vorgängerdateien mit ihren Git-Blob-
Identitäten, die exakten Schurblöcke aller 33 Historien, Nullrichtungen,
Ränge, Phasen und Konstanteneinschließungen. Die mathematische Aussage stützt
sich auf den analytischen Beweis; Softwarekontrollen allein wären kein Ersatz.

\textbf{Ausgeführte Kontrollen: 11 Testgruppen.}
Die Enumeration, Hesse-Matrizen und $F_4$-Werte werden gegen die getrennt
eingefrorenen 33 Hesse-Zertifikate geprüft. Weitere Gruppen kontrollieren alle
Innenlücken, Regulatorableitungsmajoranten, Kreuzränge, die beiden schiefen
Nullrichtungen, die drei persistierenden Historien, deren Labelmultiplizität,
beide Frequenzen, Gaußdämpfungen, den Staub-Nichtnullbeweis sowie Paket- und
Maßfaktoren. Symbolische Integrale prüfen die gewichteten Cauchy--Schwarz-
Konstanten. Ballarithmetik bestätigt (19) und die Normschranken.

Zusätzlich vergleicht eine unabhängige direkte dreidimensionale komplexe
Gaußintegration mit der gestuften Schur-Elimination für jede der 33 Historien,
zwei positive Regulatorwerte und zwei Ausgangspunkte: 132 Vergleiche.
Verwendet werden Eigenwertwurzeln statt derselben Pivotimplementierung.
Dies ist ein Float64-Strukturtest mit relativer/absoluter Toleranz
$2\cdot10^{-12}\max(1,|\text{Referenz}|)$, kein Zertifikat für einen
unendlichen Integrationsbereich. Der Nicht-$L^2$-Beweis bleibt rein analytisch.

\textbf{Reproduktion.} Python mit \texttt{sympy}, \texttt{numpy} und
\texttt{python-flint}; die konkret verwendeten Versionen stehen in
\texttt{environment.json}. Im entpackten Verzeichnis:
\begin{verbatim}
python -B test_domain.py
python -B verify.py
\end{verbatim}
Der Verifier kontrolliert alle Paketprüfsummen, neun Herkunfts-Blobs und einen
bytegleichen Ergebnis-Replay in einem temporären Verzeichnis. Das PDF wird aus
dieser vollständigen \LaTeX-Herleitung erzeugt. Ein fehlgeschlagener Kontrolllauf
ist ein Fehler und darf nicht als bestandene Reproduktion ausgewiesen werden.

\textbf{Interpretation bisheriger Resultate.} Die endlichen Matrixpaarungen aus
PR \#282--283 bleiben unverändert. Ihre Verwendung als Tests der Amplitude ist
durch (3)--(6) gerechtfertigt. Eine Interpretation als Spalte eines auf dem
bisherigen Randraum definierten Evolutionsoperators scheitert bereits für (11).
Die zuvor bedingte Bessel-Deutung einer endlichen ausgelassenen Norm ist für
diesen Eingang keine endliche globale Fehlerbilanz; die vollständige Antwort
liegt außerhalb $\HH$. Die Fehlergrößen (6)--(8) gelten stattdessen für die
kontrollierten Paarungen bzw. die $\VV'$-Antwort.

\textbf{Nächster sachlicher Anschluss.} Ein Verklebungsgesetz für solche
Amplituden benötigt passende Zwischenräume und Konvergenzbedingungen.
$K:\VV\to\VV'$ kann nicht automatisch mit sich selbst komponiert werden,
weil $\VV'$ nicht im Eingangsbereich $\VV$ liegt. Zusätzlich bleiben die bereits
gezeigten schnittübergreifenden Wirkungsterme erforderlich. Eine physikalische
Norm, ein Constraint-/Rigging-Verfahren oder zusätzliche Randdaten müssen aus
weiteren Modellannahmen begründet werden. Eine ausgehende Gaußdämpfung würde
den Operator ändern und ist keine Reparatur derselben Rechnung.

\textbf{Unverändert offen.} Physikalisches Historienmaß, vollständige
Gravitationsrandwirkung, vierdimensionale bzw. unendliche Historiengrenzen,
physikalisches Skalarprodukt, Unitarität, klassischer Grenzfall, Rückwirkung,
Horizonte, globaler Außenraum und Singularitätsauflösung. Alle acht Gates
\texttt{MODEL, CURRENT, CORRESPONDENCE, BOUNCE, REGULARITY, BOUNDARY, TRAPPING,
EXTENSION} bleiben \texttt{HOLD}. Keine Produktionssimulation, kostenpflichtige
Rechnung oder neue Messreihe wurde ausgeführt.

\section*{Quellen und Herkunft}
Die neue Definitionsbereichsanalyse und ihr Gegenbeispiel sind eigene
Ableitungen aus den hier eingefrorenen Modellformeln. Es wird keine externe
Publikation als Beleg für diesen neuen Satz ausgegeben.
\begin{itemize}
\item \href{https://github.com/thomasmgfeller/x6-spectral-astrolabium/pull/279}{PR \#279}: minimale gemeinsame Skalarwirkung.
\item \href{https://github.com/thomasmgfeller/x6-spectral-astrolabium/pull/280}{PR \#280}: Randmaß, normierte Pakete und kinematischer Hilbertraum.
\item \href{https://github.com/thomasmgfeller/x6-spectral-astrolabium/pull/281}{PR \#281}: endliche Amplitudenexistenz und Innenlücke.
\item \href{https://github.com/thomasmgfeller/x6-spectral-astrolabium/pull/282}{PR \#282}: genaue 33er-Gewichtung, gekoppelte Wirkung und Phasenkonvention.
\item \href{https://github.com/thomasmgfeller/x6-spectral-astrolabium/pull/283}{PR \#283}: Zustandsreduktion, Verklebung und bisher nur entkoppelte $L^2$-Obstruktion.
\end{itemize}
\end{document}
